\section{Introduction}
\label{sec:intro}
Modeling and forecasting realized volatility is a cornerstone of empirical finance, with direct applications in risk management, derivatives pricing, and portfolio optimization. The Heterogeneous Autoregressive (HAR) model of \cite{corsi2009simple} has become the benchmark for this task due to its parsimony and its ability to capture the long-memory properties of volatility through a simple additive cascade of daily, weekly, and monthly realized variance components.

However, financial markets are increasingly driven by exogenous information flows, macroeconomic shocks, and retail sentiment that may not be fully absorbed into the autoregressive history of prices. The proliferation of natural language processing (NLP), notably models like FinBERT \cite{araci2019finbert, yang2020finbert}, allows researchers to quantify market sentiment from news and social media at high frequencies. A burgeoning literature asks whether such alternative data adds incremental predictive power beyond the HAR baseline. Early attempts to linearly augment the HAR model with exogenous variables often suffer from the curse of dimensionality; while in-sample fit improves, out-of-sample forecasting performance deteriorates due to overfitting to noisy sentiment signals.

This paper makes four distinct contributions to this literature. First, we conduct a rigorous cross-asset evaluation spanning a 24/7 retail-driven cryptocurrency market (Bitcoin), a deeply liquid institutional equity market (S\&P 500), and an emerging market equity index (NIFTY 50). Second, we propose a hybrid architecture, \textit{Neural-HAR}, which uses a linear ordinary least squares (OLS) HAR prior and a deep residual ensemble composed of a Gated Residual Network (GRN) and an FT-Transformer \cite{gorishniy2021revisiting} to digest non-linear sentiment and microstructure interactions. Third, we perform a regime-conditional analysis using a Hidden Markov Model to test whether the marginal value of sentiment varies across market stress regimes. Fourth, we test the economic significance of these forecasts using a volatility-targeting framework.

Our empirical results present a nuanced picture. We find that linear sentiment-augmented models (HAR-S) generally underperform the HAR baseline out-of-sample. The non-linear Neural-HAR model, however, isolates predictive signal, outperforming the HAR baseline significantly under the QLIKE loss function and surviving in the Model Confidence Set \cite{hansen2011model}. Despite this statistical superiority, the economic significance is mixed: the Neural-HAR model does not improve risk-adjusted returns (Sharpe ratio) net of turnover costs, but it does significantly reduce maximum drawdowns and absolute realized volatility.
